% Generated from living-morally.md - do not edit.

\chapter{Origins of Objective Morals}

\begin{quote}
In the face of religious polemics, counting ourselves as 'subjectivists' or 'relativists' puts us at a significant social and political disadvantage. --- Leslie Allan, Director Humanists of Australia\index{Allan, Leslie}
\end{quote}

Six different theories on morality that require no divinity:

\begin{itemize}
\item Ideal Observer (David Hume)
\item Categorical Imperative (Immanuel Kant)
\item Utilitarianism (Jeremy Bentham, John Stuart Mill, Peter Singer)
\item Prescriptivism (Richard Hare)
\item Social Contract (John Rawls)
\item Stoic Cosmopolitanism (Massimo Pigliucci, John Shook)
\end{itemize}

\section{Ideal Observer (David Hume)}%
\index{Hume, David}\index{ideal observer}\nobreak%

What do we mean when we say something is good?  The idea observer theory was proposed by David Hume, Adam Smith, and Roderick Firth.  There is a hypothetical morally perfect observer.  We have to imaging this observer judging the situation.  the observer would have to be omniscient, omnipercipient, impartial, dispassionate, consistent, and otherwise normal.  Such person will be able to unlock the moral truths.  When we say X is wrong, what we are saying is that an ideal observer would judge it as wrong.  This is a kind of moral realism.\index{Smith, Adam}\index{Firth, Roderick}

The weakness of this theory is that nobody is omniscient, and it is hard for us to even know what an omniscient person would be like.  It seems like circular reasoning: if morally good is defined what an ideal observer would see, and then it is not surprising that an ideal observer would find that to be moral.\index{omniscient}

\section{Categorical Imperative (Immanuel Kant)}%
\index{Kant, Immanuel}\index{categorical imperative}\nobreak%

People have moral duties which they must follow as a duty.  Kant's theory focusses on the morality of actions, which is different than the morality of consequences.  ``Act only according to that maxim whereby you can at the same time will that it should become a universal law.''   The problem is that conditionality can not be attached to the imperatives. Each act has to be defined as an independent act.  For example lying: Kant has nothing to say about the Nazis that come to the Anne Frank door: lying is simply lying and it is better that people die than that you say a falsehood.  We know that context matters, and that an an act can not defined in such a way that it is independent of context.

\section{Utilitarianism (Jeremy Bentham, John Stuart Mill, Peter Singer)}%
\index{Bentham, Jeremy}\index{Mill, John Stuart}\index{Singer, Peter}\index{utilitarianism}\nobreak%

Utility is that aspect of an action that corresponds to whether the action augments or diminished the happiness of the party in question.  Bentham ties the moral value of an action to the amount of happiness that it produces, and/or the amount of pain it reduces.  Everything is a means to something else, until you get to pleasure.  Pleasure is an end in itself.  Pain is an end in itself as well.  Aggregate across all people to find out the total good.

Utilitarianism is a kind of individualistic version of consequentialism.  the only thing that matters is the result of your actions.  If you do something that nobody even notices, then there is no consequence of the action, and the action is neither good nor bad.

This is an interesting theory because: it provides an independent test for moral teaching.  It gives us a straightforward test to apply to moral claims.  It renders a verdict on any possible action.  The three subtheses it rests on are all plausible.

\section{Prescriptivism (Richard Hare)}%
\index{Hare, Richard}\index{prescriptivism}\nobreak%

Moral words are prescriptive.  Moral judgements are actually orders to do or not do something.

\section{Social Contract (John Rawls)}%
\index{Rawls, John}\index{social contract}\nobreak%

\section{Stoic Cosmopolitanism (Massimo Pigliucci, John Shook)}%
\index{Pigliucci, Massimo}\index{Shook, John}\index{stoic cosmopolitanism}\nobreak%
